% ============================================================
% Appendix --- Mid-semester repetition quiz
% Originally: content/week-08.tex, \section{Repetition quiz} (Corsin Week 8).
% It reviews the whole differential-topology half of the course rather than any one
% topic, so under a topic structure it belongs to no chapter and is collected here.
% ============================================================
\chapter{Mid-Semester Repetition Quiz}
\label{ch:repetition_quiz}

% Transition: Claude Opus 5 (Medium)
Corsin devoted the first half of one exercise class to a quiz revisiting everything up to this
point: differentiability, the chain rule, extrema, the implicit function theorem,
submanifolds. It ranges across \crefrange{ch:prerequisites}{ch:submanifolds} and belongs to no
single chapter, so it is collected here, where it also serves as a self-test before the
integration half of the course.

% Source: Corsin Nick/Class Notes/Week 8.pdf, p. 2
\section{The quiz}
\label{sec:repetition_quiz}

\begin{ainote}
Corsin credits the questions to Prof.\ Serra and Prof.\ Lang, and ran them from a shared Kahoot
at the mid-semester repetition session. On his pages, \textcolor{green!60!black}{green} marks the
\emph{correct} statement and \textcolor{red!80!black}{red} an \emph{incorrect} one; that colouring
is rendered below as explicit \textbf{True}/\textbf{False} verdicts.
\end{ainote}

\begin{question}
\label{q:quiz_q1}
Let $U := \{(x,y) \in \mathbb{R}^2 : 1 \leq x^2+y^2 \leq 4\}$ and
$f(x,y) := \sin(xy) - y^4$. Which of the following holds?
\begin{enumerate}[label=\textbf{(\alph*)}]
  \item $f(U) = [a,b]$ for some $a < b$.
  \item $f(U)$ is disconnected and closed with two connected components.
\end{enumerate}
\end{question}

\begin{answer}
\textbf{(a) is true, (b) is false.} $U$ is compact and connected and $f$ is continuous, so
$f(U) \subseteq \mathbb{R}$ is compact (\cref{item:continuity_preserves_compact}) and connected
(\cref{thm:continuity_preserves_connected}). The connected subsets of $\mathbb{R}$ are exactly
the intervals, and a compact one is closed and bounded, so $f(U) = [a,b]$.
\end{answer}

\begin{ainote}
Worth noting how little computation this needs: $f$ is never evaluated. The annulus $U$ is
connected (it is \emph{not} simply connected, but that is a different property and irrelevant
here), and the two preservation theorems do the rest. Recognising that a question
is asking about preserved properties rather than about values is most of the work.
\end{ainote}

\begin{question}
\label{q:quiz_q2}
Let $X \subseteq \mathbb{R}^2$ with $X \neq \mathbb{R}^2$ and $X \neq \emptyset$. If $X$ is
complete, is it closed?
\end{question}

\begin{answer}
\textbf{True.} Let $(a_n)_{n=0}^\infty$ be a convergent sequence with elements in $X$. Then
$(a_n)$ is Cauchy, and by completeness it has a limit point in $X$. So $X$ is sequentially
closed, hence closed (\cref{prop:open_closed_via_sequences}).
\end{answer}

\begin{ainote}
The limit is taken in $\mathbb{R}^2$, and completeness of $X$ is what forces it back into $X$.
The converse also holds here, but only because the ambient space $\mathbb{R}^2$ is itself
complete: a closed subset of a complete space is complete. In a non-complete ambient space
closedness gives nothing. Note that $\mathbb{Q}$ is closed in $\mathbb{Q}$ and not complete.
\end{ainote}

% Source: Corsin Nick/Class Notes/Week 8.pdf, p. 3
\begin{question}
\label{q:quiz_q3}
Let $X \subseteq \mathbb{R}^2$ as above. If $X$ is not open, is it closed?
\end{question}

\begin{answer}
\textbf{False.}
\end{answer}

% Generator: Claude Opus 5 (Medium) --- Corsin gives the verdict without a counterexample.
\begin{ainote}
Corsin marks this false but gives no counterexample; the half-open square
$X := [0,1] \times (0,1)$ from \cref{sec:open_closed_sets} is one, being neither open
nor closed. \qt{Open} and \qt{closed} are not opposites and not exhaustive: $\emptyset$ and
$\mathbb{R}^2$ are both, most sets are neither, and the excluded values
$X \neq \emptyset,\ X \neq \mathbb{R}^2$ in the hypothesis are there precisely to rule out the
two clopen cases.
\end{ainote}

\begin{question}
\label{q:quiz_q4}
Let
\[f(x,y,z) := \frac{xy^2}{x^2+y^2+z^2} \quad \text{for } (x,y,z) \neq 0, \qquad f(0,0,0) := 0.\]
Is $f \in C^1(\mathbb{R}^3)$?
\end{question}

\begin{answer}
\textbf{False.} We compute
\[\partial_xf = \frac{y^2}{x^2+y^2+z^2} - 2\,\frac{x^2y^2}{(x^2+y^2+z^2)^2}.\]
In spherical coordinates $x = r\sin\theta\cos\varphi$, $y = r\sin\theta\sin\varphi$,
$z = r\cos\theta$, every power of $r$ cancels:
\[\partial_xf = \sin^2\theta\sin^2\varphi - 2\sin^4\theta\cos^2\varphi\sin^2\varphi.\]
The limit $\lim_{r\to0}\partial_xf(r,\theta,\varphi)$ is therefore ill-defined, since it depends
on the direction $(\theta,\varphi)$ along which the origin is approached. Consequently
$\partial_xf$ is not sequentially continuous at $0$.
\end{answer}

\begin{ainote}
The cancellation of $r$ is the whole argument, and it is worth naming the reason: $\partial_xf$
is \newterm{homogeneous of degree $0$}, so it is constant along each ray from the origin and
takes a different constant on different rays. Compare
\Cref{ex:all_directional_derivatives_exist}, where the same trick (restrict to a ray, watch the
parameter cancel) exposed a function whose directional derivatives all existed without it being
differentiable.
\end{ainote}

% Source: Corsin Nick/Class Notes/Week 8.pdf, pp. 3--4
\begin{question}
\label{q:quiz_q5}
Let $f \in C^\infty(\mathbb{R}^3)$ and $\alpha \in \mathbb{R}$ be such that
\[\nabla f(0) = (0,0,0), \qquad
\Hess f(0) = \begin{pmatrix} 1 & \alpha & 0 \\ \alpha & 2 & 0 \\ 0 & 0 & 1\end{pmatrix}.\]
Which of the following is \textbf{false}?
\begin{enumerate}[label=\textbf{(\alph*)}]
  \item For all $\alpha \in \mathbb{R}$, $0$ is not a local maximum point of $f$.
  \item For $\alpha > \sqrt2$, $0$ is a saddle point.
  \item For $\alpha \neq \pm\sqrt2$, $\nabla f : \mathbb{R}^3 \to \mathbb{R}^3$ is a local
    diffeomorphism around $0$.
  \item For $\alpha < \sqrt2$, $0$ is a local minimum point of $f$.
\end{enumerate}
\end{question}

\begin{answer}
\textbf{(d) is the false one.}

\textbf{(a) is true,} since $1 > 0$ is clearly an eigenvalue of $\Hess f(0)$ for any choice of
$\alpha$ (the third standard basis vector is an eigenvector), so $\Hess f(0)$ is never negative
semidefinite and $0$ can never be a local maximum.

\textbf{(b) is true.} Let $\lambda_1,\lambda_2,\lambda_3 = 1$ be the eigenvalues of
$\Hess f(0)$. Then
\[\lambda_1\lambda_2 = \det\big(\Hess f(0)\big) = 2-\alpha^2 < 0 \quad \text{for } \alpha > \sqrt2,\]
so $\Hess f(0)$ has both positive and negative eigenvalues, making it indefinite, and $0$ is a
saddle point by \cref{thm:hessian_test}.

\textbf{(c) is true.} For these values $\det(\Hess f(0)) \neq 0$, and therefore the inverse
function theorem (\cref{thm:inverse_function_theorem}) applies. Note that
$\Hess f(0) = D(\nabla f)_0$: the Hessian of $f$ \emph{is} the differential of the gradient map.

\textbf{(d) is false,} because the same reasoning as in \textbf{(b)} gives that for
$\alpha < -\sqrt2 < \sqrt2$ we again have $2-\alpha^2 < 0$, so $0$ is a saddle point of $f$, not
a local minimum. The statement is true only on $-\sqrt2 < \alpha < \sqrt2$.
\end{answer}

\begin{ainote}
Corsin's third bullet reads \qt{for $\alpha \neq \sqrt2$}, but $\det(\Hess f(0)) = 2-\alpha^2$
vanishes at $\alpha = -\sqrt2$ as well, so it must read $\alpha \neq \pm\sqrt2$; typeset with
$\pm$ above. The oversight is a conspicuous one here, since the \emph{next} bullet is false for
exactly the reason the third one omits. The quiz is built around the missing negative root.

\end{ainote}

% Source: Corsin Nick/Class Notes/Week 8.pdf, p. 5
\begin{question}
\label{q:quiz_q6}
Let $V := \{(x,y) \in \mathbb{R}^2 : y^4+y^2 = x^3+x\}$. Decide:
\begin{enumerate}[label=\textbf{(\alph*)}]
  \item $V$ is a graph over $y$ in a neighbourhood of $(0,0)$.
  \item $V$ is a graph over $x$ in a neighbourhood of $(0,0)$.
\end{enumerate}
\end{question}

\begin{answer}
\textbf{(a) True.} Notice that $V = f^{-1}\{0\}$ for $f(x,y) := y^4+y^2-x^3-x$. Then
\[\Jac f(x,y) = (\partial_xf,\ \partial_yf) = (-3x^2-1,\ 4y^3+2y), \qquad \Jac f(0,0) = (-1,\ 0).\]
Since $\partial_xf(0,0) = -1 \neq 0$, the implicit function theorem
(\cref{thm:implicit_function_theorem}) applies with $x$ as the variable to solve for: in a
neighbourhood of $(0,0)$,
\[(x,y) \in V \iff f(x,y) = 0 \iff x = x(y).\]
The theorem does \textbf{not} apply for $y$, since $\partial_yf(0,0) = 0$.

\textbf{(b) False.} Solving the defining equation for $y^2$ by the quadratic formula gives
\[y^2 = -\tfrac12 + \sqrt{\tfrac14 + x^3 + x},\]
which has no solution for $x \in (-\varepsilon,0)$: there $x^3+x < 0$, so the square root is
smaller than $\tfrac12$ and the right-hand side is negative, while $y^2 \geq 0$. So $V$ contains
no points at all just to the left of the origin, and cannot be a graph over $x$ there.
\end{answer}

\begin{ainote}
The two parts are settled by different means, and the asymmetry is the point. Part \textbf{(a)}
is decided by \emph{checking a derivative}; part \textbf{(b)} needs an \emph{explicit
computation} showing the set is empty on one side. That is because the implicit function theorem
yields no negative results: failure of its hypothesis proves nothing on its own, exactly as
discussed after \cref{thm:implicit_function_theorem}.
\end{ainote}

% Source: Corsin Nick/Class Notes/Week 8.pdf, p. 6
\begin{question}
\label{q:quiz_q7}
Let $\gamma : I \to \mathbb{R}^2$ be injective and smooth, with $I$ an open interval. Is
$\gamma(I)$ then a smooth submanifold of $\mathbb{R}^2$?
\end{question}

\begin{answer}
\textbf{False.} Consider
\[\gamma : (0,2\pi) \to \mathbb{R}^2, \qquad t \mapsto \big(\sin t,\ \sin 2t\big),\]
whose image is a \newterm{figure eight} \germanfor{Lemniskate}. This $\gamma$ is injective and
smooth, but there is no graphical representation of $\operatorname{im}(\gamma)$ around $(0,0)$.
\end{answer}

% Sonnet 5 (Medium) --- FIG-W08-01, redrawn from the plot on the source page.
\begin{center}
\begin{tikzpicture}[>=Latex, scale=1.45]
  \draw[->] (-1.35,0) -- (1.35,0) node[right, font=\scriptsize] {$x$};
  \draw[->] (0,-1.25) -- (0,1.25) node[above, font=\scriptsize] {$y$};
  \draw[blue!80!black, thick, variable=\t, domain=0:6.2832, samples=200, smooth]
    plot ({sin(\t r)}, {sin(2*\t r)});
  \draw[red!80!black, thick, dashed] (0,0) circle (0.33);
  \node[red!80!black, font=\scriptsize, right] at (0.52,0.34) {problem!};
  \draw[red!80!black, thin, ->] (0.50,0.30) -- (0.20,0.13);
\end{tikzpicture}
\end{center}

\begin{ainote}
Two things are worth extracting. First, how injectivity survives a self-crossing at all: the two
parameter ends $t \to 0^+$ and $t \to 2\pi^-$ both send $\gamma(t) \to (0,0)$, but neither
\emph{attains} it, so the image passes through the origin only once (via $t = \pi$) while
being approached from two further directions.

Second, this locates precisely which hypothesis of \cref{thm:local_parametrization} fails. It is
not injectivity, and not injectivity of $D\gamma_t$; it is the requirement that the
parametrization be a \textbf{homeomorphism} onto its image. Here $\gamma^{-1}$ is discontinuous
at $(0,0)$: points of the image arbitrarily close to the origin have parameters near $0$ and near
$2\pi$, far apart in $I$. That the resulting set is not a submanifold can then be confirmed by
the component count of \cref{ex:cross_not_submanifold}. Deleting the crossing point leaves
four pieces, while deleting a point from an interval leaves two.
\end{ainote}

% Sonnet 5 (Medium)
The quiz ends here, and with it the differential-topology half of the course. The rest of the
week starts integration in several variables.
